% Topic 4: Using Matrix Algebra for Systems of Linear Equations in Two Variables
% Part of the Matrix Fundamentals section

\subsection{Using Matrix Algebra for Systems of Linear Equations in Two Variables}

Any system of linear equations can be translated into a single matrix equation. This allows us to use the power of matrix inverses to solve for the variables directly, bypassing manual substitution or elimination.

\subsubsection*{Forming the Matrix Equation}

Consider the general $2 \times 2$ system:
\[
\begin{cases}
ax + by = e \\
cx + dy = f
\end{cases}
\]
By using the rules of matrix multiplication, we can separate the coefficients, the variables, and the constants into three matrices:
\[
\begin{bmatrix} a & b \\ c & d \end{bmatrix} \begin{bmatrix} x \\ y \end{bmatrix} = \begin{bmatrix} e \\ f \end{bmatrix}
\]
This is often written compactly as $AX = B$, where:
\begin{itemize}[leftmargin=*]
  \item $A$ is the \textbf{coefficient matrix}.
  \item $X$ is the \textbf{variable matrix} (a column vector).
  \item $B$ is the \textbf{constant matrix} (a column vector).
\end{itemize}

\subsubsection*{Solving the Matrix Equation}

To isolate the variable matrix $X$, we cannot ``divide'' by matrix $A$. Instead, we multiply both sides of the equation by the inverse matrix, $A^{-1}$.

Because matrix multiplication is not commutative, we must multiply from the \textbf{left} on both sides:
\begin{align*}
AX &= B \\
A^{-1} (AX) &= A^{-1} B \\
(A^{-1} A) X &= A^{-1} B \\
I X &= A^{-1} B \\
X &= A^{-1} B
\end{align*}

If $\det(A) = 0$, then $A^{-1}$ does not exist. This algebraically confirms the situations discussed in the previous topic: the system has either no unique solution (parallel lines) or infinitely many solutions (coincident lines).

\bigskip

% ── Worked Example: Solving a System using Matrices ──
\begin{problem}[Matrix Solution to a $2 \times 2$ System]
Solve the following system of equations using matrix algebra:
\begin{align*}
2x + 3y &= 8 \\
5x - y &= 3
\end{align*}
\emph{(Note: This is the same system solved algebraically in Topic 3.)}
\end{problem}

\begin{solution}
\textbf{Step 1: Write the system as a matrix equation.} \\
Ensure the equations are in standard form (variables on the left, constants on the right). They are. \\
$AX = B \implies \begin{bmatrix} 2 & 3 \\ 5 & -1 \end{bmatrix} \begin{bmatrix} x \\ y \end{bmatrix} = \begin{bmatrix} 8 \\ 3 \end{bmatrix}$

\medskip
\textbf{Step 2: Find the determinant of $A$.} \\
$\det(A) = (2)(-1) - (3)(5) = -2 - 15 = -17$. \\
Since $\det(A) \neq 0$, a unique solution exists.

\medskip
\textbf{Step 3: Find the inverse matrix, $A^{-1}$.} \\
$A^{-1} = \frac{1}{-17} \begin{bmatrix} -1 & -3 \\ -5 & 2 \end{bmatrix} = -\frac{1}{17} \begin{bmatrix} -1 & -3 \\ -5 & 2 \end{bmatrix}$

\medskip
\textbf{Step 4: Solve for $X$ using $X = A^{-1}B$.} \\
$X = -\frac{1}{17} \begin{bmatrix} -1 & -3 \\ -5 & 2 \end{bmatrix} \begin{bmatrix} 8 \\ 3 \end{bmatrix}$ \\[6pt]
Calculate the matrix product first: \\
$X = -\frac{1}{17} \begin{bmatrix} (-1)(8) + (-3)(3) \\ (-5)(8) + (2)(3) \end{bmatrix}$ \\[6pt]
$X = -\frac{1}{17} \begin{bmatrix} -8 - 9 \\ -40 + 6 \end{bmatrix} = -\frac{1}{17} \begin{bmatrix} -17 \\ -34 \end{bmatrix}$ \\[6pt]
Finally, multiply the scalar into the matrix: \\
$X = \begin{bmatrix} 1 \\ 2 \end{bmatrix}$

\medskip
\textbf{Conclusion:} \\
Since $X = \begin{bmatrix} x \\ y \end{bmatrix} = \begin{bmatrix} 1 \\ 2 \end{bmatrix}$, the solution is $x = 1, y = 2$. \\
This precisely matches the geometric point of intersection $(1, 2)$ found previously!
\end{solution}

\begin{takeaways}
\item \textbf{Left-Multiplication is Critical:} You must write $X = A^{-1}B$. Writing $X = BA^{-1}$ is mathematically incorrect because the inner dimensions for multiplication would not match (a $2 \times 1$ matrix cannot be multiplied by a $2 \times 2$ matrix in that order).
\item \textbf{Scalar Timing:} As seen in the example, delay multiplying the scalar fraction (like $-\frac{1}{17}$) until the very last step. This avoids messy fraction arithmetic during the matrix multiplication.
\item \textbf{Determinant Check:} The determinant provides an immediate check. If $\det(A) = 0$, you can stop and state that no unique geometric intersection exists.
\end{takeaways}
